\documentclass[12pt]{article}
\usepackage{amsmath,amsfonts}
\usepackage{fullpage}
\usepackage{hyperref}
\usepackage[utf8]{inputenc}

\begin{document}

TP2, ROP631, été 2026, Prof. Toby HOCKING, à rendre le 10 juillet.

\section{exercice 3.2.1}
Soit $h(x+d\theta)$ la fonction pour recherche linéaire, à partir d’un point $x\in \mathbb R^n$, dans la direction $d\in\mathbb R^n$, d’un longueur (pas) $\theta>0$.
Considérons le problème
$\min_{a\le\theta\le b}h(\theta)$. Si un point $\theta_0$ est
dans l'intérieur de l'intervalle $[a,b]$, alors le théorème 2.2.1
s'applique directement; lorsque $\theta_0=a$ ou $\theta_0=b$, il
suffit que $h$ soit croissante (en $a$) ou décroissante (en $b$)
pour avoir un minimum local sous contrainte: en effet, on ne compare
$a$ qu'avec des valeurs $\theta\ge a$. Obtenez une condition
nécessaire et suffisante pour que $a$ ou $b$ soient des minima
locaux de $h$. L'exercice consiste donc à fournir les conditions encadrées et démontrer:
\begin{quotation}
  $\theta^*$ est min local de la fonction $h$ si et seulement si
  \begin{itemize}
  \item[$\bullet$] $\theta^*=a$ et \fbox{conditions pour que $h$ soit croissante}   ou
  \item[$\bullet$] $a<\theta^*<b$ et la première dérivée non nulle de $h$ est d'ordre pair et de signe positif ou
  \item[$\bullet$] $\theta^*=b$ et \fbox{conditions pour que $h$ soit décroissante.}
  \end{itemize}
\end{quotation}

\section{exercice 3.2.7}
Dans cet exercice, nous explorons un problème d'ajuste\-ment de
fonction. Les valeurs suivantes de $x$ et $y$ sont données: 
$$\begin{tabular}{|l||l|l|l|l|}
\hline
x&0&1&2&3\\
\hline
y&0&0&3&9\\
\hline
\end{tabular}$$
On soupçonne qu'une relation d'ordre 3 relie $x$ à $y$, et on envisage
un modèle du type $y=a_0 + a_3x^3$.
  \begin{enumerate}
    \item Formulez le problème 
    d'identifier les valeurs de $a_0$ et $a_3$ qui minimisent la somme des carrés
    des différences entre le modèle et les données.
    \item Exprimez les conditions d'optimalité (nécessaires et
    suffisantes, de premier et second ordre) du problème en a).
    \item Nous allons traiter une simplification du problème en supposant que $a_0=0$, donc que la relation recherchée
    est purement cubique (nouveau problème avec seulement $a_3$). On analyse ici le problème de dimension un qui ne comporte que la variable $a_3$. Cette
    simplification semble raisonnable puisque pour $x=0$, $y=0$.
    \begin{enumerate}
    \item calculez la valeur de $a_3$ qui minimise la somme des carrés
    des différences entre le modèle simplifié et les données; 
    \item est-ce que les conditions suffisantes d'optimalité sont
    satisfaites pour la solution du problème de dimension un en i)?
    \end{enumerate}
    \item Revenons maintenant au modèle de dimension deux en a) pour calculer les valeurs de $a_0$ et $a_3$. 
    Calculez les valeurs de $a_0$ et $a_3$ du modèle
     avec l'algorithme du gradient conjugué.
Est-ce que les conditions suffisantes d'optimalité sont
    satisfaites?  
\end{enumerate}

\section{exercice 3.3.1}
Pour la fonction $x^2+2y^2+4x+4y$, vérifiez par induction que l'algorithme du gradient amorcé à l'origine engendre la suite $\begin{pmatrix}{2\over3^k}-2\\ \left(-{1\over3}\right)^k-1\end{pmatrix}$. Déduisez-en le minimum de la fonction.

\section{exercice 3.11.2}

Utilisez Julia pour implanter deux algorithmes d’optimisation : comparez les résultats de l'algorithme d'Armijo (avec direction = négative gradient), avec l’algorithme de régions de confiance (avec pas de Newton calculé avec le gradient conjugué).
Votre programme utilisera trois fonctions qui retournent respectivement $f(x)$, $\nabla f(x)$ et $\nabla^2f(x)$.
Testez vos programmes avec (1) la fonction que nous avons étudié dans les démos, $f(x)=(x_1^2 + x_2 - 11)^2 + (x_1+x_2^2-7)^2$, et (2) la fonction de Rosenbrock : $f(x)=100(x_2-x_1^2)^2 + (1-x_1)^2$ à partir du point de départ $(-1.2,1.0)^t$. Cette fonction a longtemps constitué un défi pour les programmes de minimisation. Les modifications de l'algorithme de Newton traitent habituellement cette fonction facilement.

Comme dans les démos, faire un graphique (ou tableau) avec les valeurs de $f(x_k)$ et $||\nabla f(x_k)||$ au fil des itérations $k$.

Est-ce que vous voyez des vitesses de convergence linéaire / quadratique comme attendu ?

Est-ce que vos résultats sont consistents avec les démos ?

Vous pouvez éventuellement vous inspirer des codes sources (dans R) pour les démos,
\url{https://github.com/tdhock/2026-05-nlopt}

\end{document}
